\documentclass[12pt]{article}

\usepackage[utf8]{inputenc}
\usepackage[romanian]{babel}
\usepackage{amsmath, amssymb, amsthm}
\usepackage{geometry}
\geometry{a4paper, margin=2.5cm}

\title{Ipostaze ale Intuiției în Predarea Matematicii}
\author{}
\date{}

\theoremstyle{definition}
\newtheorem{definition}{Definiție}

\theoremstyle{plain}
\newtheorem{proposition}{Propoziție}

\begin{document}

\maketitle

\begin{proposition}
Ipostaze ale intuiției în transmiterea și interiorizarea cunoștințelor de matematică:
\begin{enumerate}
    \item Exprimările \emph{de ordin intuitiv} sau \emph{în limbaj cotidian} permit o înțelegere mai rapidă și o memorare mai lesnicioasă a afirmațiilor matematice.
    \item Suportul intuitiv conduce la un model care contribuie substanțial la rezolvarea problemei.
    \item Intuiția sugerează metoda de rezolvare.
    \item Intuiția poate furniza răspunsuri la întrebarea problemei.
    \item Uneori, intuiția poate furniza o justificare convingătoare care însă nu constituie o demonstrație.
    \item Anumite abordări intuitive, în special cele de ordin geometric, pot conduce la demonstrații complete.
\end{enumerate}
\end{proposition}

\end{document}
